\documentclass[UTF8,11pt]{ctexart}
\usepackage[a4paper,margin=24mm]{geometry}
\usepackage{amsmath,amssymb,amsthm,mathtools}
\usepackage[all]{xy}
\usepackage{xurl,hyperref,enumitem}
\hypersetup{hidelinks,breaklinks=true}
\setlength{\emergencystretch}{3em}
\newtheorem*{theorem}{Theorem}
\newtheorem*{lemma}{Lemma}
\newtheorem*{proposition}{Proposition}
\newtheorem*{corollary}{Corollary}
\theoremstyle{definition}
\newtheorem*{definition}{Definition}
\newtheorem*{remark}{Remark}
\DeclareMathOperator{\Hom}{Hom}
\DeclareMathOperator{\Ext}{Ext}
\DeclareMathOperator{\Tor}{Tor}
\DeclareMathOperator{\Spec}{Spec}
\DeclareMathOperator{\supp}{supp}
\DeclareMathOperator{\im}{im}
\DeclareMathOperator{\id}{id}
\DeclareMathOperator{\Tr}{Tr}
\DeclareMathOperator{\rank}{rank}
\newcommand{\R}{\mathbb R}
\newcommand{\C}{\mathbb C}
\newcommand{\Z}{\mathbb Z}
\newcommand{\N}{\mathbb N}
\newcommand{\Q}{\mathbb Q}
\begin{document}
\section*{019\quad 球面连续映射的Borsuk--Ulam重合点定理}
\subsection*{来源与收录范围}
Haynes Miller, \emph{Algebraic Topology I}, MIT 18.905, Fall 2016. 根据 Sanath Devalapurkar 的课堂记录编写并经授课者修订的讲义；原书另将图示归于 Xianglong Ni。主命题为 Theorem 38.11 (Borsuk--Ulam)，印刷第107页（合订 PDF 第109页）。采用 OCW 提供的 Fall 2016 文件，获取日期 2026-09-28。
\par\url{https://ocw.mit.edu/courses/18-905-algebraic-topology-i-fall-2016/2dad3dd8354992d877fbeca4049817ee_MIT18_905F16_lecture_notes.pdf}

\subsection*{作者命题与人类证明}

\begin{theorem}[38.11, Borsuk--Ulam]
Think of $S^n$ as the unit vectors in $\R^{n+1}$. For any continuous function $f:S^n\to\R^n$, there exists $x\in S^n$ such that $f(x)=f(-x)$.
\end{theorem}
\begin{proof}
Suppose that no such $x$ exists. Then we may define a continuous function $g:S^n\to S^{n-1}$ by
\[
g:x\longmapsto\frac{f(x)-f(-x)}{\|f(x)-f(-x)\|}.
\]
Note that $g(-x)=-g(x)$: $g$ is equivariant with respect to the antipodal action. It descends to a map $\bar g:\mathbf{RP}^n\to\mathbf{RP}^{n-1}$.

We claim that $\bar g_*:H_1(\mathbf{RP}^n)\to H_1(\mathbf{RP}^{n-1})$ is nontrivial. To see this, pick a basepoint $b\in S^n$ and choose a $1$-simplex $\sigma:\Delta^1\to S^n$ such that $\sigma(e_0)=b$ and $\sigma(e_1)=-b$. The group $H_1(\mathbf{RP}^n)$ is generated by the cycle $p\sigma$. The image of this cycle in $H_1(\mathbf{RP}^{n-1})$ is represented by the loop $\bar g p\sigma$ at $\bar b=pb$, which is the image of the $1$-simplex $g\sigma$ joining $gb$ to $g(-b)=-g(b)$. The class of this $1$-simplex thus generates $H_1(\mathbf{RP}^{n-1})$.

Therefore $\bar g$ is nontrivial in $H_1(-;\mathbb F_2)$, and hence also in $H^1(-;\mathbb F_2)$. Writing $a_n$ for the generator of $H^1(\mathbf{RP}^n;\mathbb F_2)$, we must have $a_n=\bar g^*a_{n-1}$, and consequently
\[
a_n^n=(\bar g^*a_{n-1})^n=\bar g^*(a_{n-1}^n).
\]
But $H^n(\mathbf{RP}^{n-1};\mathbb F_2)=0$, so $a_{n-1}^n=0$; while $a_n^n\neq0$. This is a contradiction.
\end{proof}

\subsection*{整理说明（非作者正文）}
只计 Borsuk--Ulam 一个问题，不把构造等变球面映射或其射影空间映射另计题数。必要先修是覆盖空间的路径提升、模二同调与上同调的对应及射影空间的上同调环
\[
H^*(\mathbf{RP}^m;\mathbb F_2)=\mathbb F_2[a]/(a^{m+1}),\qquad |a|=1.
\]
此环公式在作者紧邻的 Example 38.10 明示；本题承担从假设构造等变映射、证明其在一次同调上的非平凡作用，再用乘法结构产生矛盾，并非把 Borsuk--Ulam 或“不存在奇映射”作为先修。

正文由实际取得的 PDF 文本转录。将球面映射与商映射分别排作 $g$ 与 $\bar g$，商基点排作 $\bar b$；$p$ 为源文中的射影商映射。原文的基点措辞、未另列低维情形的写法均保留，不额外补造证明。该页截图接口未成功，未声称逐页图像校对，也未保存下载 PDF 原件。本条是经典代数拓扑构造性证明，难度依据等变性如何转成杯积障碍，不按定理名称或篇幅判断。
\subsection*{可修改的简短标注（非作者正文）}
$c$：球面连续映射的反足重合点。

$a$：归一化差映射；反足群作用；射影空间商；覆盖空间路径；一次同调；模二上同调环；杯积的自然性。

\texttt{cross\_theory\_status = yes}。把一个连续映射的重合点问题转成等变球面映射，再转成射影空间上非平凡的一次上同调拉回；杯积幂与维数消失迫使矛盾。位置：Theorem 38.11 的完整证明。

全部标注均为非穷尽、可修改初稿；未标注不表示未使用。
\subsection*{许可与修改记录}
原文来自 MIT OpenCourseWare，作者与课程如上。按来源网站标示的 Creative Commons Attribution--NonCommercial--ShareAlike 4.0 许可复用；本转录及整理沿用同一许可。2026-09-28：助手转录、加入来源定位与简短标注；不暗示作者或 MIT 认可本次整理。许可：\url{https://creativecommons.org/licenses/by-nc-sa/4.0/}。
\end{document}
